\documentclass{article}
\usepackage[utf8]{inputenc}
\usepackage{amsmath,amssymb}
\usepackage[most]{tcolorbox}
\usepackage{geometry}
\geometry{margin=2.5cm}

\newcommand{\step}[1]{\par\noindent\hangindent=3.5em\hangafter=1%
  \makebox[3.5em][l]{\textbf{#1.}}}
\newcommand{\steptag}[1]{\hfill{\small\textsf{[#1]}}}

\newtcolorbox{proofbox}[1]{
  colback=gray!5, colframe=gray!60,
  fonttitle=\bfseries, title={#1},
  breakable, enhanced,
  left=4pt, right=4pt, top=4pt, bottom=4pt,
  before skip=10pt, after skip=10pt
}

\begin{document}

\begin{proofbox}{Quantitative inverse-function control: if V:X->Y is a Ban...}
\step{1} Quantitative inverse-function control: if V:X->Y is a Banach-space isomorphism and f:B\_r(x\_0)->Y is C\textasciicircum{}1 with ||V\textasciicircum{}(-1)Df(x)-I|| <= c < 1, then f is injective, (1-c)||x\_1-x\_2|| <= ||V\textasciicircum{}(-1)(f(x\_1)-f(x\_2))|| <= (1+c)||x\_1-x\_2||, and f(B\_r(x\_0)) contains f(x\_0)+V(B\_\{(1-c)r\}(0)). \steptag{validated} \\[6pt]
\step{1.1} Secant control and injectivity: under the root hypotheses, for all x\_1,x\_2 in B\_r(x\_0), (1-c)||x\_1-x\_2|| <= ||V\textasciicircum{}(-1)(f(x\_1)-f(x\_2))|| <= (1+c)||x\_1-x\_2||; because 1-c>0, f is injective. \steptag{validated} \\[6pt]
\step{1.1.1} Define F:B\_r(x\_0)->X by F(x)=V\textasciicircum{}(-1)(f(x)-f(x\_0)). Then F is C\textasciicircum{}1, DF(x)=V\textasciicircum{}(-1)Df(x), F(x\_0)=0, and the root hypothesis is exactly ||DF(x)-I\_X||<=c throughout B\_r(x\_0). \steptag{validated} \\[4pt]
\step{1.1.2} Secant-error lemma: if a C\textasciicircum{}1 map F on the open ball B\_r(x\_0) satisfies ||DF(x)-I\_X||<=c throughout that ball, then for all x\_1,x\_2 in the ball, ||F(x\_1)-F(x\_2)-(x\_1-x\_2)||<=c||x\_1-x\_2||. \steptag{validated} \\[6pt]
\step{1.1.2.1} Fix x\_1,x\_2 in B\_r(x\_0), put d=x\_1-x\_2 and gamma(t)=x\_2+td. Convexity of the ball keeps gamma([0,1]) in the domain. For q(t)=F(gamma(t))-F(x\_2)-td, the chain rule gives q'(t)=(DF(gamma(t))-I\_X)d and hence ||q'(t)||<=c||d|| for every t in [0,1]. \steptag{validated} \\[4pt]
\step{1.1.2.2} Banach-valued Newton-Leibniz estimate: if q:[0,1]->X is C\textasciicircum{}1 and ||q'(t)||<=M, then q(1)-q(0) is the norm-limit of the derivative Riemann sums (equivalently the Bochner/Riemann integral of q'), so ||q(1)-q(0)||<=integral\_0\textasciicircum{}1 ||q'(t)||dt<=M. Applying this with M=c||d|| and observing q(1)-q(0)=F(x\_1)-F(x\_2)-d proves the secant-error estimate. \steptag{validated} \\[6pt]
\step{1.1.2.2.1} For any C\textasciicircum{}1 curve q:[0,1]->X with ||q′(t)||<=M, the Banach-valued fundamental theorem of calculus gives q(1)-q(0)=integral\_0\textasciicircum{}1 q′(t)dt: indeed, continuity of q′ makes its tagged Riemann sums converge in norm to that integral, while the corresponding sums of increments telescope to q(1)-q(0). The norm inequality for the Bochner/Riemann integral then gives ||q(1)-q(0)||<=integral\_0\textasciicircum{}1||q′(t)||dt<=M. Now fix x\_1,x\_2 in B\_r(x\_0), put d=x\_1-x\_2, gamma(t)=x\_2+td, and q(t)=F(gamma(t))-F(x\_2)-td. Validated node 1.1.2.1 establishes that gamma([0,1]) lies in the domain and that q is C\textasciicircum{}1 with ||q′(t)||<=c||d||. Apply the preceding estimate with M=c||d||. Since gamma(1)=x\_1 and gamma(0)=x\_2, q(1)=F(x\_1)-F(x\_2)-d and q(0)=0, hence ||F(x\_1)-F(x\_2)-(x\_1-x\_2)||=||q(1)-q(0)||<=c||d||=c||x\_1-x\_2||. \steptag{validated} \\[4pt]
\step{1.1.3} Norm deduction: for vectors a,d in a normed space, ||a-d||<=c||d|| implies (1-c)||d||<=||a||<=(1+c)||d|| by the reverse and ordinary triangle inequalities; when c<1, a=0 forces d=0. \steptag{validated} \\[4pt]
\step{1.2} Quantitative image inclusion: under the root hypotheses, every y in f(x\_0)+V(B\_\{(1-c)r\}(0)) equals f(x) for some x in B\_r(x\_0). \steptag{validated} \\[6pt]
\step{1.2.1} Let g:B\_r(0)->X be g(h)=V\textasciicircum{}(-1)(f(x\_0+h)-f(x\_0)). Then g(0)=0 and ||Dg(h)-I\_X||<=c. Applying the segment argument of the secant-error lemma gives ||(g(h)-g(k))-(h-k)||<=c||h-k|| for every h,k in B\_r(0). \steptag{validated} \\[4pt]
\step{1.2.2} Centered covering with the correct specialization: any map g:B\_r(0)->X with g(0)=0 and ||(g(h)-g(k))-(h-k)||<=c||h-k|| for 0<=c<1 satisfies B\_\{(1-c)r\}(0) subset g(B\_r(0)). For the image-inclusion application, g is specifically the map defined in node 1.2.1, g(h)=V\textasciicircum{}(-1)(f(x\_0+h)-f(x\_0)); for this g, if y=f(x\_0)+Vz and g(h)=z, then y=f(x\_0+h). \steptag{validated} \\[6pt]
\step{1.2.2.1} Fix z with ||z||<(1-c)r. Choose rho with 0<rho<r and ||z||<(1-c)rho, let K=\{h in X:||h||<=rho\}, and define T(h)=z+h-g(h). For h,k in K the secant-error hypothesis gives ||T(h)-T(k)||<=c||h-k||. Taking k=0 and g(0)=0 gives ||h-g(h)||<=c||h||, hence ||T(h)||<=||z||+c rho<rho; thus T maps the complete closed ball K into itself and is a contraction. \steptag{validated} \\[4pt]
\step{1.2.2.2} Fixed-point construction and translation: let (M,d) be a complete metric space and let T:M->M be c-Lipschitz with 0<=c<1. For h\_\{n+1\}=T(h\_n), induction gives d(h\_\{n+1\},h\_n)<=c\textasciicircum{}n d(h\_1,h\_0). Thus, for m>n, the triangle inequality gives d(h\_m,h\_n)<=sum\_\{j=n\}\textasciicircum{}\{m-1\} c\textasciicircum{}j d(h\_1,h\_0)<=c\textasciicircum{}n d(h\_1,h\_0)/(1-c), so (h\_n) is Cauchy and completeness gives h\_n->h in M. Moreover d(T(h),h)<=d(T(h),T(h\_n))+d(h\_\{n+1\},h)<=c d(h,h\_n)+d(h\_\{n+1\},h)->0, hence T(h)=h. Apply this to the contraction T on the complete closed ball K from 1.2.2.1, with its norm-induced metric. Then h lies in K subset B\_r(0), and T(h)=h means z+h-g(h)=h, equivalently g(h)=z. Therefore every z in B\_\{(1-c)r\}(0) lies in g(B\_r(0)); and if y=f(x\_0)+Vz, then g(h)=z implies f(x\_0+h)=f(x\_0)+Vg(h)=y, with x\_0+h in B\_r(x\_0). \steptag{validated} \\[4pt]
\step{1.2.2.3} Specialization and translation, explicitly using node 1.2.1: take g(h)=V\textasciicircum{}(-1)(f(x\_0+h)-f(x\_0)), rather than an arbitrary map satisfying only the secant bound. For any z in B\_\{(1-c)r\}(0), the generic covering construction in nodes 1.2.2.1 and 1.2.2.2 supplies h in B\_r(0) with g(h)=z. The defining identity for this particular g is Vg(h)=f(x\_0+h)-f(x\_0). Hence g(h)=z implies f(x\_0+h)-f(x\_0)=Vz, and therefore for y=f(x\_0)+Vz one has y=f(x\_0+h); also x\_0+h lies in B\_r(x\_0). No such translation is asserted for an independently chosen admissible g. \steptag{validated} \\[4pt]
\end{proofbox}

\end{document}
